% =========================================================
% Ordinary Differential Equations --- Breadcrumb
% =========================================================

\begin{tcolorbox}[
  colback=gray!6,
  colframe=gray!40,
  arc=2pt,
  left=8pt, right=8pt, top=6pt, bottom=6pt,
  title={\small\textbf{Where This Chapter Lives}},
  fonttitle=\small\bfseries
]
\begin{center}
\small
Calculus
$\;\to\;$ \textbf{Ordinary Differential Equations}
$\;\to\;$ Dynamical Systems / Numerical Methods / PDE
\end{center}

\medskip
\noindent\textbf{Why this chapter belongs here.}
Ordinary differential equations are where derivatives stop being only local
rates and start governing whole families of evolving states.  Once
differentiation is in hand, equations such as
\[
y'(t)=f(t,y(t))
\]
become the natural language for motion, growth, decay, forcing, and
feedback.

\medskip
\noindent\textbf{What this chapter will gather.}
The eventual ODE chapter is the right home for:
\begin{itemize}
  \item first-order equations and separable models;
  \item linear equations and integrating factors;
  \item existence and uniqueness ideas;
  \item systems of ODEs and phase portraits;
  \item stability, equilibrium, and qualitative behavior;
  \item numerical methods such as Euler and Runge--Kutta schemes.
\end{itemize}

\medskip
\noindent\textbf{Why this is currently a breadcrumb.}
The surrounding curriculum is still building the supporting architecture:
calculus, linear algebra, and numerical analysis.  ODEs sit at
the meeting point of all three, so this file marks the destination clearly
even before the chapter is fully written.
\end{tcolorbox}
